% Ad Familiares 2.13 --- headnote
% ad-familiares-02-13, early May 50 BC, Laodicea

\textit{Cicero to M. Caelius Rufus, curule aedile,
written from Laodicea between the Kalends and the Nones
of May 50 BC (Perseus dateline: \textit{Scr. Laudiceae
inter K. et Non. Maias a. 704 (50)}) --- so within the
first week of May, in the closing weeks of Cicero's
Cilician assizes and a day or two before the date
(Nones, 7 May) he names in section 3 as his intended
departure for the summer camp. A reply to a Caelius
newsletter whose rarity Cicero notes in the opening
line: the post from Rome is barely getting through, but
when it does Caelius's reports are as shrewd and as
welcome as ever.}

\textit{Three movements. Cicero defends his standing
with Appius Claudius Pulcher, his predecessor in
Cilicia: there is no breach, only a difference in the
shape of their administrations, and Cicero has now made
himself an advocate against the danger Appius is in
from Dolabella's recent prosecution. He then turns to
the Roman political weather Caelius has sketched ---
the ``lethargy of the state'' and the astonishing news
that Curio, of all people, has come over and is now
defending Caesar. The closing sentences open out into
the homesick refrain that will recur in every letter
of this Cilician spring and summer: the assizes are
done, the books are clean, the allies are unaggrieved,
and Cicero wants only to set the soldier in his summer
quarters and come home to see Caelius hold his
aedileship.}
